% !TEX root = ../main.tex
\documentclass[../main.tex]{subfiles}

\begin{document}

\selectlanguage{vietnamese}

\begin{center}
    \textbf{\LARGE TÀI LIỆU ÔN THI TOÀN DIỆN CƠ SỞ TOÁN HỌC 5 CHƯƠNG AI} \\
    \vspace{0.2cm}
    \textbf{\large Hệ thống hóa Đặc trưng Thuật toán qua Biểu thức Toán học theo Đề cương Học phần} \\
    \textit{(Biên soạn chuẩn xác theo cấu trúc Đề cương môn học và Giáo trình Russell \& Norvig 3rd Edition)}
\end{center}

\vspace{0.5cm}

\section{CHƯƠNG 1: TỔNG QUAN VỀ TRÍ TUỆ NHÂN TẠO (Introduction)}
\subsection{1.1. Định nghĩa và khái niệm Trí tuệ nhân tạo (Artificial intelligence definition and concepts)}
Định nghĩa AI theo giáo trình Russell \& Norvig được cấu trúc hóa dựa trên ma trận hai chiều: Hệ thống suy nghĩ (\textit{Thought processes}) đối lập với Hành vi thực tế (\textit{Behavior}), và Mô phỏng con người (\textit{Human performance}) đối lập với Tính hợp lý tối ưu (\textit{Rationality}).
\begin{itemize}
    \item \textit{Thinking Humanly:} Mô hình hóa nhận thức của bộ não thông qua các mô phỏng máy tính và thực nghiệm tâm lý học.
    \item \textit{Acting Humanly:} Đạt được năng lực hành động như con người thông qua việc vượt qua Phép thử Turing (\textit{Turing Test}) - yêu cầu tích hợp các nhánh xử lý ngôn ngữ tự nhiên, biểu diễn tri thức, suy luận tự động, và học máy.
    \item \textit{Thinking Rationally:} Sử dụng cấu trúc tam đoạn luận và các quy luật logic (\textit{Laws of Thought}) để thiết lập các lập luận đúng đắn tuyệt đối.
    \item \textit{Acting Rationally (Trọng tâm giáo trình):} Thiết kế các đại lý thông minh hành động hợp lý (\textit{Rational Agent}) để tối đa hóa hàm hiệu suất kỳ vọng trong môi trường.
\end{itemize}

\subsection{1.2. Lịch sử của Trí tuệ nhân tạo (History of AI)}
Sự tiến hóa của AI trải qua các giai đoạn nền tảng: Giai đoạn ấp ủ (1943--1955) với mô hình nơ-ron số đầu tiên của McCulloch \& Pitts; Sự ra đời của thuật ngữ AI tại hội thảo Dartmouth (1956); Giai đoạn nhiệt huyết và kỳ vọng lớn (1952--1969); Giai đoạn hiện thực phũ phàng và mùa đông AI thứ nhất (1966--1973); Sự bùng nổ của các Hệ chuyên gia dựa trên tri thức (1969--1979); Mùa đông AI thứ hai (1984--1987); và Kỷ nguyên hiện đại của AI dựa trên dữ liệu lớn, toán học xác suất thống kê và học sâu từ năm 1995 đến nay.

\subsection{1.3. Các lĩnh vực nghiên cứu và ứng dụng chính (Main research and application areas)}
Các nhánh chuyên sâu cốt lõi bao gồm: Thị giác máy tính (\textit{Computer Vision}), Xử lý ngôn ngữ tự nhiên (\textit{NLP}), Hệ thống tự hành và Robot (\textit{Robotics}), Quy hoạch chiến lược tự động (\textit{Automated Planning}), Hệ quản trị tri thức, và Học máy (\textit{Machine Learning}).

\subsection{1.4. Thành tựu và những bài toán chưa có lời giải (Achievements and unsolved problems)}
\begin{itemize}
    \item \textit{Thành tựu:} Vượt qua con người trong các trò chơi chiến lược phức tạp (Cờ vua Deep Blue, Cờ vây AlphaGo), xe tự lái cấp độ cao, hệ thống chẩn đoán y tế tự động chuẩn xác và các mô hình ngôn ngữ lớn (LLM).
    \item \textit{Bài toán chưa giải quyết:} Thách thức về Trí tuệ nhân tạo tổng hợp (\textit{AGI} - Trí tuệ nhân tạo đa dụng cấp chuyên gia), tính minh bạch và khả năng giải thích của mô hình hộp đen (\textit{Explainable AI}), hiện tượng ảo giác dữ liệu, và giới hạn tính toán lý thuyết đối với các bài toán có không gian trạng thái bùng nổ hàm mũ (NP-hard).
\end{itemize}

\end{document}
